\chapter{What drives the performance}
\label{ch:drivers}
\Cref{ch:hybrid} showed that \hybrid{} scores 0.11 above \tfunet{} on LOFAR and
matches the best published results. The natural explanation is that the
components added for RFI---the strip convolutions, channel attention and
residual blocks---are responsible. This chapter tests that explanation
directly. It turns out to be wrong.

\section{Removing the components one at a time}
\label{sec:ablation}
To find out what each component contributes, I wrote a version of the network
in which each component can be switched off, and trained every variant
through exactly the same pipeline as the full model: same data, splits,
normalisation, loss, optimiser, number of steps and threshold selection. Only
the network changes. As a check, the switchable network with everything
switched on scored 0.6624, in line with the real model's $0.6603 \pm 0.0040$, so
it faithfully reproduces \hybrid{}.

\Cref{tab:ablation} shows the results. Removing the strip convolutions,
the attention or the residual blocks one at a time changes the score by less
than 0.007, and removing the strip convolutions even raised it slightly.
Removing all three at once gives 0.6605, essentially the same as the full
model. We call this model the ``partial hybrid'': it is \hybrid{} without its
three added modules, and it is not the original U-Net or \tfunet{}, since it
keeps all the other differences listed in \cref{tab:hybrid_vs_tfunet}.

\begin{table}[H]
  \centering
  \caption[Ablation of the network components (one seed)]{Ablation on the 109 LOFAR test images (maximum \Fone{}, one seed).
    The seed-to-seed standard deviation of the full model is 0.0040, so
    differences below about 0.008 are within noise.}
  \label{tab:ablation}
  \begin{tabular}{lrrr}
    \toprule
    Variant & Parameters & max \Fone{} & Change \\
    \midrule
    Full model                                   & 593{,}842 & 0.6624 & --- \\
    without strip convolutions                   & 508{,}210 & 0.6637 & $+0.0013$ \\
    without channel attention                    & 593{,}818 & 0.6557 & $-0.0067$ \\
    without residual blocks                      & 572{,}074 & 0.6569 & $-0.0055$ \\
    without all three (partial hybrid)           & 486{,}418 & 0.6605 & $-0.0019$ \\
    partial hybrid without GroupNorm             & 485{,}682 & 0.6289 & $-0.0335$ \\
    \bottomrule
  \end{tabular}
\end{table}

With one run per variant, only the last row looked like a real effect:
removing GroupNorm cost 0.034. Because a single run can mislead, I repeated the two most
important variants with two more seeds (\cref{tab:ablation3}).

\begin{table}[H]
  \centering
  \caption[Ablation over three seeds]{The key ablation variants over three seeds (maximum \Fone{}).}
  \label{tab:ablation3}
  \small
  \begin{tabular}{lcccc}
    \toprule
    Model & Seed 0 & Seed 1 & Seed 2 & Mean $\pm$ s.d. \\
    \midrule
    Full \hybrid{}                 & 0.6592 & 0.6647 & 0.6570 & $0.6603 \pm 0.0040$ \\
    Partial hybrid (three components removed) & 0.6605 & 0.6618 & 0.6531 & $0.6585 \pm 0.0047$ \\
    Partial hybrid without GroupNorm  & 0.6289 & 0.6600 & 0.6595 & $0.6495 \pm 0.0178$ \\
    \bottomrule
  \end{tabular}
\end{table}

The full model and the partial hybrid are not significantly different (Welch
$t$-test, $p = 0.64$): the three components contribute nothing measurable. The
apparent effect of GroupNorm also largely disappeared. The first run without it
was an outlier; the other two scored the same as the partial hybrid, and with
three seeds the difference is not significant ($p = 0.48$). The one clear sign
is that without GroupNorm the results vary about four times more from seed to
seed, so normalisation seems to make training more stable even if it does not
make the final model better on average.

\section{Adding a CRF refinement}
\label{sec:crf}
A conditional random field (CRF) is an extra step after the network that
changes each pixel's label to agree better with its neighbours. Chen et
al.\footcite{chen2023} used one to find weak RFI in pulsar data, so I added one
on top of \hybrid{}. It did not help. With Chen et al.'s settings the score fell
from 0.6592 to 0.4884, because the CRF deleted the thin, scattered RFI that is
common in LOFAR; with the best settings it rose by only 0.0012, too little to
count. The network's convolutions already look at each pixel's neighbours, so
the CRF had nothing new to add.

\section{What explains the gain over \tfunet{}?}
\label{sec:unexplained}
Taken together, these results show that the gain of \hybrid{} over \tfunet{}
does not come from the components it was designed around. The partial hybrid,
with these components removed but the same training setup, scores almost as
well. What remains are the other
differences listed in \cref{tab:hybrid_vs_tfunet}: the Dice loss, the output
layer without a ReLU, same padding, the different depth and width, batch size,
learning rate and framework.

\subsection{The Dice loss}
\label{sec:dice}
The first of these I tested is the loss. I trained \hybrid{} again with
exactly the same settings and seeds, but with plain cross-entropy, the loss
\tfunet{} uses, instead of cross-entropy plus Dice. Since the components were
shown above to make no measurable difference, the full model was used, so that
each run pairs with a headline run of the same seed.

\begin{table}[H]
  \centering
  \caption[Effect of the Dice loss]{\hybrid{} with and without the Dice loss on the 109 LOFAR test
    images (mean $\pm$ standard deviation over three seeds).}
  \label{tab:dice}
  \small
  \begin{tabular}{lrrrr}
    \toprule
    Loss & max \Fone{} & pooled \Fone{} & ROC AUC & PR AUC \\
    \midrule
    with Dice (headline) & $0.6603 \pm 0.0040$ & $0.6561 \pm 0.0044$ & $0.938 \pm 0.011$ & $0.658 \pm 0.009$ \\
    without Dice         & $0.6495 \pm 0.0049$ & $0.6472 \pm 0.0064$ & $0.987 \pm 0.001$ & $0.708 \pm 0.008$ \\
    \tfunet{}            & $0.5482 \pm 0.0139$ & $0.5021 \pm 0.0259$ & $0.947 \pm 0.007$ & --- \\
    \bottomrule
  \end{tabular}
\end{table}

Without Dice the maximum \Fone{} fell for all three seeds, by 0.007 to 0.015
(mean 0.011; paired $t = -4.44$ with 2 degrees of freedom, significant at the
5\% level), and the pooled \Fone{} by 0.009 (\cref{tab:dice}). The Dice loss
is therefore a real part of the gain, but a small one: it accounts for 0.011
of the 0.112 by which \hybrid{} leads \tfunet{}, about a tenth. Even without
it, \hybrid{} scores 0.10 above \tfunet{}.

As with class weighting (\cref{ch:hybrid}), the ROC area moved the other way:
without Dice it rose from 0.938 to 0.987, and the area under the
precision--recall curve from 0.658 to 0.708. The Dice loss improves the
decision at the best threshold but makes the model rank pixels less well
overall, probably because it pushes outputs towards exactly 0 or 1. It also
explains why the headline model's ROC area is slightly below \tfunet{}'s.

The remaining 0.10 lies among the output layer, padding, depth and width,
batch size, learning rate and framework. I have not tested these one by one,
so I cannot yet say which of them is responsible.
